\section{Introduction}
Feldman, Kamath, and Srebro ask whether distribution-independent statistical-query learning forces a common low-dimensional linear representation \citep{FKS26}. We answer their SQ question negatively, including its polynomial and probabilistic-dimension variants. The distinction is between \emph{distribution independence}, which asks for one algorithm valid under every marginal, and \emph{distribution adaptation}, which permits that algorithm to ask where the unknown marginal places mass.

A linear representation must fix every example--hypothesis sign at once. A learner can instead query marginal masses, label regions on which its surviving hypotheses agree, and use disagreements to eliminate hypotheses. A product-distribution rectangle condition is exactly what a marginal-free mixture can cover pointwise, by finite minimax. The new step is to turn that mixture into an adversarial-oracle learner whose progress is the product of the two rectangle masses. A bounded-drift argument gives a worst-case query budget, not just finite expected termination. The random incidence flips of \citet{HHPTZ22} have large rectangles and large sign-rank, so these two resources separate even for one algorithm serving all marginals.

\paragraph{Contributions.}
\begin{itemize}
\item If $\rect(H)\ge1/\rho$, RECTIFY uses $O(\rho(\log|H|+\log(1/\epsilon)))$ queries at tolerance $\Omega(\epsilon/(\rho\log(1/\epsilon)))$ to achieve expected error at most $\epsilon$ against every valid adaptive oracle. Applied to the incidence flips on $2M$ points, it refutes all three polynomial dimension implications, including $\dc_{1/2}$ and $\dc^{C\epsilon}$.
\item A fixed-dimensional strict template with no negative $2\times2$ admits a pointwise cap sampler and a finite-bit implementation polynomial in the supplied table. A finite net also gives a deterministic learner. For the flip family a proper completion has $O(\log M+1/\epsilon+\log(1/\epsilon))$ queries at the same tolerance order.
\item We prove the all-product bound $\rect(\operatorname{PG}(2,q))=\Theta(1/q)$, exhibit SQ-easy cube halfspaces with $\rect\le e^{-n/4}$, quantify necessary marginal adaptation under a small-query budget, and show that source-description length alone does not repair the implication when computation is uncharged.
\end{itemize}

\paragraph{Imported results and attribution.}
The high-sign-rank construction is due to \citet{HHPTZ22}, and its counting uses Warren's strict-sign-pattern theorem \citep{Warren68,AMY16}. The homogeneous-rectangle theorem \citep{APP05} is used only for the sharper constant $2^{-15}$; the self-contained cap proof already gives a constant sufficient for the separation. The minimax rectangle-mixture device also appears in the average-margin proof of \citet[Theorem~3.1]{HHPTZ22}; we do not claim that device as new. The new learning argument is its adaptive peeling/elimination analysis, its uniform query-budget proof, and the proper completion. The probabilistic lower bounds use the actual target-dependent marginal quantifiers, rather than uniform approximate sign-rank. Feldman's characterization is invoked only for the separate comparison in Section~6, not for the main learner. We additionally use real-closed-field decidability in the program-length counterexample and state that imported algorithmic fact explicitly. The weakly polynomial convex-program implementation needed for caps is proved in Appendix~\ref{app:convex}.

\subsection{Related work}
Forster's spectral approach connects sign-rank with normalized vector representations \citep{Forster02}. Hatami et al. compare VC dimension, average margin and monochromatic-rectangle ratio, and construct matrices that separate these parameters from sign-rank \citep{HHPTZ22}. Their geometric and counting ingredients motivate the present learner; they do not themselves provide its unknown-marginal oracle analysis.

Feldman's 2017 characterization concerns general randomized SQ search complexity, and his finite-field line example separates distribution-specific from distribution-independent SQ learning \citep{Feldman17}. Our class lies on the easy side of distribution-independent SQ learning, so that older separation does not settle the question addressed here. Approximate radial-isotropy algorithms have been developed independently, including strongly polynomial methods \citep{DKT22}. We use their geometric object, but only claim weakly polynomial computation and include a direct fixed-dimensional implementation proof. The engineered differentiable-learning simulations of \citet{AKMSS21} do not preserve the Gaussian-initialized bias-free architecture in the remaining SGD question. Recent topological sign-rank bounds and upper bounds for Hamming-distance matrices are discussed in Section~\ref{sec:resources} \citep{FHV26,GHIS25}.
